\chapter{Patterns: lazy state}
\label{ch:lazystate}

\begin{aside}
A signal that exists only when needed. Pay the cost on first
access; nothing before.
\end{aside}

\section{Definition}

\begin{lstlisting}
auto big = lazy_signal([&] {
    return parse_huge_file();
});
\end{lstlisting}

The lambda runs on first \code{get}.

\section{Caching}

Once computed, the result caches. Subsequent reads are
fast.

\section{Invalidation}

To recompute, reset the lazy:

\begin{lstlisting}
big.invalidate();
auto v = ctx.get(big);  // recomputes
\end{lstlisting}

\section{Use cases}

\begin{itemize}[leftmargin=*]
  \item Settings panel content that few users open.
  \item Help screen that loads markdown.
  \item Big derived tables.
  \item Expensive validators that depend on infrequent
    state.
\end{itemize}

\section{Startup impact}

Moving work out of startup shortens the time to first
frame. Lazy state is a powerful tool for this.

\section{Recap}

\begin{recap}
\begin{itemize}[leftmargin=*]
  \item Lazy signal computes on first read.
  \item Caches; invalidate to recompute.
  \item Shorten startup by deferring.
\end{itemize}
\end{recap}

\section*{Workshop}

\begin{enumerate}[leftmargin=*]
  \item Identify three signals you could make lazy.
  \item Convert them; measure startup time.
\end{enumerate}
